\documentclass[11pt]{article}
\usepackage[a4paper, margin=2.5cm]{geometry}

\usepackage[utf8]{inputenc}
\usepackage[T1]{fontenc}
\usepackage[slovene]{babel}
\usepackage{lmodern}

% Paketi za pomoč pri matematiki.
\usepackage{amsthm}
\usepackage{amssymb}

\usepackage{graphicx}

\usepackage{makeidx}



{\theoremstyle{plain}
\newtheorem{izrek}{Izrek}
}

{\theoremstyle{definition}
\newtheorem{definicija}{Definicija}
}

\newcommand{\f}{\mathcal{F}}

% Brownovo gibanje
% Matej Rojec

\begin{document}

\title{Brownovo gibanje}
\author{Matej Rojec}
\date{}

\maketitle


Brownovo gibanje (več v \cite{karatzas1991brownian}) je intuitivno slučajen proces, % Sklic na knjigo
ki predstavlja naključno gibanje delcev v mediju.
 
    % Slika: PerrinPlot2.pdf, velikost scale=0.3
    \begin{figure}[h]
    \centering
    \includegraphics[width=0.4\textwidth]{PerrinPlot2.pdf}
    \caption{Reprodukcija slike iz Jean Baptiste Perrin, \emph{Mouvement brownien et réalité moléculaire}, Ann. de Chimie et de Physique (VIII) 18, 5-114, 1909}
    \label{fig:perrin}
    \end{figure}
    % Napis pod sliko: 
    % Reprodukcija slike iz Jean Baptiste Perrin, \emph{Mouvement brownien et réalité moléculaire}, Ann. de Chimie et de Physique (VIII) 18, 5-114, 1909

    \begin{definicija}
    \emph{Standardno Brownovo gibanje} \(\{B_t\}_{t \geq 0}\) je slučajen proces z naslednjimi lastnostmi:
    
        
        \begin{enumerate}
            \item \(B_0 = 0\).
            \item Prirastki \(B_{t_n} - B_{t_{n-1}}, B_{t_{n-1}} - B_{t_{n-2}}, \ldots, B_{t_2} - B_{t_1}, B_{t_1} - B_{t_0}\) so neodvisne slučajne spremenljivke, za vsak \(t_0 \leq t_1 \leq \cdots \leq t_{n-1} \leq t_n\).
            \item Za vsak \(t \geq 0\) in \(h > 0\) velja \(B_{t+h} - B_t \sim \mathcal{N}(0, h)\).
            \item Funkcija \(t \mapsto B_t\) je zvezna skoraj gotovo.
        \end{enumerate}
    \end{definicija}

    
    Preden zapišemo izrek, definirajmo še pojem časa ustavljanja.
    
    \begin{definicija}
    Slučajna spremenljivka \(\tau\) na verjetnostnem prostoru \((\Omega,\f, P)\) z vrednostmi v !!
    je \emph{čas ustavljanja} glede na filtracijo \((\f_t)_{t\in{T}}\), če velja: \(\forall t\in{T}:\{r \leq t\}\in \f_t\).
    \end{definicija}
    
    Zdaj lahko zapišemo izrek \ref{thm:stopped_brownian}. % Sklic na izrek z oznako thm:stopped_brownian
    
    \begin{izrek}
    \label{thm:stopped_brownian}
    Naj bo \(\{B_t\}_{t \geq 0}\) (standardno) Brownovo gibanje, \(\tau\) čas ustavljanja glede na
    \((\f_t)_{t \geq 0}\) in naj velja \(P{[\tau \le \infty]}\).
    Potem je tudi proces:
    \[
    \hat{B} := \{B_{\tau+t} - B_\tau \mid t \geq 0\}
    \]
    (standardno) Brownovo gibanje in neodvisen od $\f_t$.
    \end{izrek}

    \bibliographystyle{plain}

    \bibliography{knjiga}

    \printindex
    
\end{document}